
\begin{figure}[htbp]
\centering
\resizebox{0.9\textwidth}{!}{%
\begin{tikzpicture}[font=\scriptsize, node distance=18mm]
  \tikzset{sta/.style={draw, circle, fill=lightaccent, minimum size=9mm, align=center},
           ap/.style={draw, rounded corners, fill=warnbg, minimum width=13mm, minimum height=8mm, align=center}}
  \node[sta] (a) {A};
  \node[ap, right=of a] (b) {Vevő\\AP};
  \node[sta, right=of b] (c) {C};
  \draw[thick] (a) -- node[above]{hallják egymást} (b);
  \draw[thick] (b) -- node[above]{hallják egymást} (c);
  \draw[dashed, red] (a) .. controls +(0,-1.2) and +(0,-1.2) .. node[below]{A és C nem hallják egymást} (c);

  \node[draw, rounded corners, fill=black!5, align=left, below=12mm of b, text width=72mm] (steps)
    {1. A RTS-t küld a vevőnek.\\2. A vevő CTS-t küld, amit C is hall.\\3. A és C a Duration mezők alapján NAV-ot állít.\\4. Az adatkeret védettebb, mert C nem kezd adásba.};
  \draw[arrow] (b) -- (steps);
\end{tikzpicture}%
}
\caption{Rejtett állomás probléma és RTS/CTS alapötlet}
\label{fig:zv23_rts_cts}
\end{figure}
